\subsubsection{边值问题的 Green 函数}\label{subsec:green-functions-boundary-problems}

考虑边值问题 \eqref{eq:general-homogeneous-boundary-value-problem}, 并采用上一节的假设.

\begin{Definition}
\label{def:green-function-boundary-problem}
\eqref{eq:general-homogeneous-boundary-value-problem} 的 \textbf{Green 函数} $G(\bm{x},\bm{y})$ 是下列问题的解
\begin{equation}\label{eq:green-function-boundary-value-definition}
\begin{aligned}
L(\bm{x},D_{\bm{x}})G(\bm{x},\bm{y})&=\delta(\bm{x}-\bm{y}),\quad\bm{x},\bm{y}\in\Omega,\\
B_j(\bm{x},D_{\bm{x}})G(\bm{x},\bm{y})&=0,\quad\bm{x}\in\partial\Omega,\ \bm{y}\in\Omega,\ j=1,\ldots,p.
\end{aligned}
\end{equation}

\end{Definition}

\eqref{eq:green-function-boundary-value-definition} 中的第一个方程应在分布意义下理解. 当然, 要使边界条件有意义, 对 $G$ 的光滑性要求必须高于仅仅是一个分布. 但对椭圆边值问题, 事实上只要 $\bm{x}\ne\bm{y}$, $G$ 就光滑, 因而解释边界条件并无困难. 显然, Green 函数的概念推广了基本解的概念. 若 $L$ 为常系数, 将 Green 函数视为基本解的扰动往往是有利的. 具体地, 若 $\mathcal{G}(\bm{x}-\bm{y})$ 为基本解, 令 $G(\bm{x},\bm{y})=\mathcal{G}(\bm{x}-\bm{y})+g(\bm{x},\bm{y})$, 其中 $g$ 满足
\begin{equation}\label{eq:green-fundamental-solution-smooth-perturbation}
L(\bm{x},D_{\bm{x}})g(\bm{x},\bm{y})=0,\quad\bm{x},\bm{y}\in\Omega,
\end{equation}
并且对 $j=1,\ldots,p$ 满足
\begin{equation}\label{eq:green-perturbation-boundary-data}
B_j(\bm{x},D_{\bm{x}})g(\bm{x},\bm{y})=-B_j(\bm{x},D_{\bm{x}})\mathcal{G}(\bm{x}-\bm{y}),\quad\bm{x}\in\partial\Omega,\ \bm{y}\in\Omega.
\end{equation}
若 $\mathcal{G}$ 在 $\bm{x}\ne\bm{y}$ 时光滑, 则 \eqref{eq:green-perturbation-boundary-data} 的右端对每个 $\bm{y}\in\Omega$ 都光滑. 对椭圆问题, 将在 Chapter~\ref{chap:linear-elliptic-equations} 中看到, 这意味着 $g$ 也光滑. 在 $\Omega$ 内部, 基本解贡献了奇性最强的部分, 因而在某种意义上 ``支配'' Green 函数.

当然, 找出哪些边值问题具有存在且唯一的 Green 函数, 是根本性的重要问题. 目前还没有处理这一问题所需的技术, 但将在 Chapter~\ref{chap:linear-elliptic-equations} 中讨论椭圆方程的存在性问题.

若 Green 函数存在, 则 \eqref{eq:general-homogeneous-boundary-value-problem} 的一个形式解为
\begin{equation}\label{eq:green-homogeneous-boundary-solution-representation}
u(\bm{x})=\int_\Omega G(\bm{x},\bm{y})f(\bm{y})\,\mathrm{d}\bm{y};
\end{equation}
事实上, 若 $f\in\mathcal{D}(\Omega)$, 则对 $G$ 作相当弱的假设, \eqref{eq:green-homogeneous-boundary-solution-representation} 就给出 \eqref{eq:general-homogeneous-boundary-value-problem} 的一个解. 例如, 只要 $G(\cdot,\bm{y})$ 作为 $\mathcal{D}^{\prime}(\Omega)$ 的元素连续依赖于 $\bm{y}$, 且 $G$ 在 $\bm{x}\ne\bm{y}$ 时光滑, 就已足够. 特别地, 若边值问题 \eqref{eq:general-homogeneous-boundary-value-problem} 唯一可解, 则 \eqref{eq:green-homogeneous-boundary-solution-representation} 导出恒等式
\begin{equation}\label{eq:green-test-function-representation-identity}
u(\bm{x})=\int_\Omega G(\bm{x},\bm{y})L(\bm{y},D_{\bm{y}})u(\bm{y})\,\mathrm{d}\bm{y}
\end{equation}
对所有 $u\in\mathcal{D}(\Omega)$ 成立.

现在假设 $G$ 具有足够的正则性, 使 \eqref{eq:green-homogeneous-boundary-solution-representation} 不仅对 $u\in\mathcal{D}(\Omega)$, 而且对每个 $u\in X_B$ 都成立. 利用 \eqref{eq:general-boundary-operator-green-identity}, 得到
\begin{equation}\label{eq:green-adjoint-kernel-boundary-derivation}
\begin{aligned}
u(\bm{x})&=\int_\Omega u(\bm{y})L^*(\bm{y},D_{\bm{y}})G(\bm{x},\bm{y})\,\mathrm{d}\bm{y}\\
&\quad+\int_{\partial\Omega}\sum_{j=1}^p S_j(\bm{y},D_{\bm{y}})u(\bm{y})C_j(\bm{y},D_{\bm{y}})G(\bm{x},\bm{y})\,\mathrm{d}S_{\bm{y}}.
\end{aligned}
\end{equation}
若这对任意 $u\in X_B$ 成立, 则对每个 $\bm{x}\in\Omega$, 必须有
\begin{equation}\label{eq:green-kernel-adjoint-variable-problem}
\begin{aligned}
L^*(\bm{y},D_{\bm{y}})G(\bm{x},\bm{y})&=\delta(\bm{x}-\bm{y}),\quad\bm{x},\bm{y}\in\Omega,\\
C_j(\bm{y},D_{\bm{y}})G(\bm{x},\bm{y})&=0,\quad\bm{y}\in\partial\Omega,\ \bm{x}\in\Omega,\ j=1,\ldots,p.
\end{aligned}
\end{equation}
也就是说, 固定 $\bm{x}$, 将 $G$ 看作 $\bm{y}$ 的函数时, 它满足伴随边值问题.

利用 \eqref{eq:green-kernel-adjoint-variable-problem}, 并在 \eqref{eq:general-boundary-operator-green-identity} 中令 $v(\bm{y})=G(\bm{x},\bm{y})$, 得到
\begin{equation}\label{eq:green-inhomogeneous-boundary-representation-identity}
\begin{aligned}
u(\bm{x})&=\int_\Omega G(\bm{x},\bm{y})L(\bm{y},D_{\bm{y}})u(\bm{y})\,\mathrm{d}\bm{y}\\
&\quad+\int_{\partial\Omega}\sum_{j=1}^p T_j(\bm{y},D_{\bm{y}})G(\bm{x},\bm{y})B_j(\bm{y},D_{\bm{y}})u(\bm{y})\,\mathrm{d}S_{\bm{y}}.
\end{aligned}
\end{equation}
因此, 若非齐次边值问题
\begin{equation}\label{eq:general-inhomogeneous-boundary-value-problem}
\begin{aligned}
L(\bm{x},D)u&=f(\bm{x}),\quad\bm{x}\in\Omega,\\
B_j(\bm{x},D)u&=\phi_j(\bm{x}),\quad\bm{x}\in\partial\Omega,\quad j=1,\ldots,p
\end{aligned}
\end{equation}
有解, 则该解表示为
\begin{equation}\label{eq:green-general-inhomogeneous-solution-formula}
\begin{aligned}
u(\bm{x})&=\int_\Omega G(\bm{x},\bm{y})f(\bm{y})\,\mathrm{d}\bm{y}\\
&\quad+\int_{\partial\Omega}\sum_{j=1}^p T_j(\bm{y},D_{\bm{y}})G(\bm{x},\bm{y})\phi_j(\bm{y})\,\mathrm{d}S_{\bm{y}}.
\end{aligned}
\end{equation}
需要提醒的是, 为论证导出 \eqref{eq:green-inhomogeneous-boundary-representation-identity} 的分部积分是合理的, $\bm{x}\in\Omega$ 这一条件十分重要, 因为它使 $G(\bm{x},\bm{y})$ 在 $\bm{y}\in\partial\Omega$ 时光滑. 一般而言, 当 $\bm{x}\in\partial\Omega$ 时, \eqref{eq:green-inhomogeneous-boundary-representation-identity} 并不表示 $u(\bm{x})$.

对简单区域内的一些简单方程, 可以显式给出 Green 函数. 作为例子, 考虑半径为 $R$ 的球 $B_R$ 上具有 Dirichlet 边界条件的 Laplace 方程. 此时, 可以用所谓的镜像法构造 Green 函数. 基本解 $\mathcal{G}(|\bm{x}-\bm{y}|)$ 可以看作位于点 $\bm{y}$ 的点电荷产生的势. 现在的想法是, 在反射点 $\overline{\bm{y}}=R^2\bm{y}/|\bm{y}|^2$ 处放置第二个点电荷, 使两个电荷的势在球面 $|\bm{x}|=R$ 上相互抵消. 这导出 Green 函数
\begin{equation}\label{eq:ball-dirichlet-green-image-formula}
\begin{aligned}
G(\bm{x},\bm{y})&=\mathcal{G}(|\bm{x}-\bm{y}|)-\mathcal{G}\left(\dfrac{|\bm{y}|}R|\bm{x}-\overline{\bm{y}}|\right)\\
&=\mathcal{G}\left(\sqrt{|\bm{x}|^2+|\bm{y}|^2-2\bm{x}\cdot\bm{y}}\right)\\
&\quad-\mathcal{G}\left(\sqrt{(|\bm{x}||\bm{y}|/R)^2+R^2-2\bm{x}\cdot\bm{y}}\right).
\end{aligned}
\end{equation}
若 $\bm{y}=0$, 令 $G(\bm{x},\bm{y})=\mathcal{G}(|\bm{x}|)-\mathcal{G}(R)$. 若 $|\bm{x}|=R$, 则 $G(\bm{x},\bm{y})=0$; 此外, 计算得
\begin{equation}\label{eq:ball-green-image-laplacian}
\Delta_{\bm{x}}G(\bm{x},\bm{y})=\delta(\bm{x}-\bm{y})-\dfrac{|\bm{y}|^2}{R^2}\delta(\bm{x}-\overline{\bm{y}}),
\end{equation}
当 $\bm{x}$ 和 $\bm{y}$ 限制在 $B_R$ 内时, 它等于 $\delta(\bm{x}-\bm{y})$. 因而 $G$ 确为 Dirichlet 问题的 Green 函数. 可以看到, $G$ 关于其两个变元对称, 这反映了 Dirichlet 问题的自伴随性.

Dirichlet 问题
\begin{equation}\label{eq:ball-dirichlet-poisson-boundary-problem}
\Delta u=f(\bm{x}),\ \bm{x}\in B_R,\qquad u=\phi(\bm{x}),\ \bm{x}\in\partial B_R
\end{equation}
的解由 \eqref{eq:green-general-inhomogeneous-solution-formula} 表示:
\begin{equation}\label{eq:ball-dirichlet-green-solution-representation}
u(\bm{x})=\int_{B_R}G(\bm{x},\bm{y})f(\bm{y})\,\mathrm{d}\bm{y}+\int_{\partial B_R}\dfrac{\partial}{\partial n_{\bm{y}}}G(\bm{x},\bm{y})\phi(\bm{y})\,\mathrm{d}S_{\bm{y}}.
\end{equation}
此外, 直接计算可得
\begin{equation}\label{eq:ball-green-boundary-normal-derivative}
\dfrac{\partial}{\partial n_{\bm{y}}}G(\bm{x},\bm{y})=\dfrac{R^2-|\bm{x}|^2}{\Omega_m R|\bm{x}-\bm{y}|^m}
\end{equation}
(参见 Problem~\ref{ex:verify-ball-green-normal-derivative}). 这就导出已在 Section~\ref{sec:dirichlet-existence-proof} 中遇到过的 Poisson 公式.
